\section{Embedding：让 Token ID 变成可训练的向量}
\label{intro:20-embeddings}
\phantomsection\label{concept:embedding}

上一章已经把文字转成整数 ID，但编号本身没有语义大小：
ID 4 与 ID 1 的差值并不是两个词语的含义差异。
我们要找一种表示，让\textbf{每个 ID 有自己的特征，而且这些特征能在训练时调整}。

\subsection{先试 One-hot，再看它缺什么}

在本书的五项教学词表中，ID 1 对应 One-hot 行向量
$(0,1,0,0,0)$，ID 3 对应 $(0,0,0,1,0)$。
它们可以区分身份，但不同类别之间的点积都为零。
真实模型词表可能很大，实际程序也不必真的创建这些长而稀疏的向量。

\subsection{同一份三维教学参数：五行向量}

从这里开始，\textbf{新版完整贯通实验的唯一初始 Embedding 表}
使用 $d_{\rm model}=3$，五个 ID 各对应一行：
\[
E=\begin{bmatrix}
0&0&0\\
1.2&0&-0.2\\
0.2&-0.4&0.7\\
-0.3&1.1&0.2\\
0.15&-0.2&1.25
\end{bmatrix}\in\R^{5\times3}.
\]
数值来自仓库 \texttt{scripts/teaching\_reference.py} 的
\texttt{parameters()}，是人为设计的、训练之前的参数，\textbf{不是真实语言模型的已学权重}。
其中“春天”是 ID 1 的 $(1.2,0,-0.2)$，“花”是 ID 3 的
$(-0.3,1.1,0.2)$，“开了”是 ID 4 的 $(0.15,-0.2,1.25)$。
三条向量\textbf{线性无关}，避免了原二维示例中
“开了 = 春天 + 花”的偶然构造被误解成语言规律。

把输入 $[1,3,4]$ 逐行查表：
\begin{equation}
\boxed{
X=\begin{bmatrix}
1.2&0&-0.2\\
-0.3&1.1&0.2\\
0.15&-0.2&1.25
\end{bmatrix}\in\R^{3\times3}.
}
\label{eq:running-embedding}
\end{equation}
每一行对应一个 Token 位置，每行有三个可训练特征。

% Uniform 5-by-3 teaching E, from scripts/teaching_reference.py
\par\noindent\begin{minipage}{\linewidth}
\tagstructbegin{tag=Figure,alttext={根据 Token ID 在五行三列的嵌入表中选取一行；ID 4 花对应负零点三、一点一、零点二}}
\begin{center}
\begin{tikzpicture}[>=Latex,x=11mm,y=7.5mm]
\node[font=\small\bfseries,anchor=south] at (1.5,1.36) {Embedding 表 $E$（$5\times3$）};
\foreach \rr/\aa/\bb/\cc in {0/0/0/0,1/1.2/0/-0.2,2/0.2/-0.4/0.7,3/-0.3/1.1/0.2,4/0.15/-0.2/1.25}{
\ifnum\rr=3 \def\bg{linegray!25}\else\def\bg{paper}\fi
\filldraw[fill=\bg,draw=linegray] (0,{-\rr}) rectangle (1,{1-\rr});
\filldraw[fill=\bg,draw=linegray] (1,{-\rr}) rectangle (2,{1-\rr});
\filldraw[fill=\bg,draw=linegray] (2,{-\rr}) rectangle (3,{1-\rr});
\node[font=\scriptsize] at (.5,{.5-\rr}) {$\aa$};
\node[font=\scriptsize] at (1.5,{.5-\rr}) {$\bb$};
\node[font=\scriptsize] at (2.5,{.5-\rr}) {$\cc$};
\node[font=\footnotesize,anchor=east] at (-.25,{.5-\rr}) {ID \the\numexpr\rr+1\relax};
}
\draw[flowcolor,very thick] (0,-3) rectangle (3,-2);
\draw[->,flowcolor,very thick] (3.16,-2.5)--(4.08,-2.5);
\node[font=\small,align=left,anchor=west] at (4.22,-2.5)
 {ID 4（“花”）\\行向量 $(-0.3,1.1,0.2)$};
\node[font=\scriptsize,anchor=west] at (0,-4.65)
 {编号从 1 开始；查表不是把 ID 本身当成特征。};
\end{tikzpicture}
\end{center}
\begin{center}\captionof{figure}{三维 Embedding 按 ID 查表；数值均为教学设定}\label{fig:embedding-table}\end{center}
\tagstructend
\end{minipage}\par


\subsection{按 ID 查表为什么等价于矩阵乘法？}

ID 3 的 One-hot 行向量右乘 $E$，只取出其中一行：
\[
\underbrace{(0,0,0,1,0)}_{1\times5}\,
\underbrace{E}_{5\times3}
=\underbrace{(-0.3,1.1,0.2)}_{1\times3}.
\]
实际程序用整数索引直接读取，不会为了查询一次 Token 就构造整条 One-hot。
训练时，这张 Embedding 表与后续网络一样，可以根据预测误差更新。

\subsection{训练前后，向量真的变了什么？}

我们先看结果，等 Transformer 部分再逐层算出它的来历。
用一层三维、单头、Pre-Norm 教学网络，以
$[1,3,4]$ 预测三个位置的目标 $[3,4,2]$，
所有参数同时用学习率 $0.04$ 做 \textbf{20 步 SGD}。
该训练样本的平均交叉熵从约 $1.7984$ 降为 $0.4832$。
这时对应的三行 Embedding 变为：
\[
\begin{aligned}
E_{20}[1]&\approx(1.2145,\,0.1364,\,-0.2272),\\
E_{20}[3]&\approx(-0.4298,\,0.8810,\,0.3345),\\
E_{20}[4]&\approx(0.1734,\,-0.2188,\,1.3136).
\end{aligned}
\]
这份结果可由 \texttt{python3 scripts/teaching\_reference.py --json}
复现；\textbf{训练损失下降只说明在这份教学题上优化成功}，不能证明模型理解真实中文。

% Real coordinates: 20 whole-network SGD steps in scripts/teaching_reference.py.
% Compare first two coordinates, not PCA and not a semantic interpretation.
\tagstructbegin{tag=Figure,alttext={三维词向量前两个坐标训练前后对照。春天从一点二和零变成一点二一四五和零点一三六四；花从负零点三和一点一变为负零点四三和零点八八；开了从零点一五和负零点二变成零点一七三四和负零点二一八八}}
\begin{center}
\begin{tikzpicture}[x=16mm,y=16mm,>=Latex]
\begin{scope}[shift={(0,0)}]
\draw[->,linegray,thick] (-.68,-.55)--(1.62,-.55) node[below,font=\scriptsize] {$e_1$};
\draw[->,linegray,thick] (-.68,-.55)--(-.68,1.46) node[above,font=\scriptsize] {$e_2$};
\node[font=\small\bfseries] at (.45,1.58) {训练前 $E_0$};
\fill[blue] (1.2,0) circle (2.5pt); \node[font=\scriptsize,above] at (1.18,.02) {春天};
\fill[teal] (-.3,1.1) circle (2.5pt); \node[font=\scriptsize,right] at (-.3,1.1) {花};
\fill[orange] (.15,-.2) circle (2.5pt); \node[font=\scriptsize,right] at (.15,-.2) {开了};
\end{scope}
\begin{scope}[shift={(3.35,0)}]
\draw[->,linegray,thick] (-.68,-.55)--(1.62,-.55) node[below,font=\scriptsize] {$e_1$};
\draw[->,linegray,thick] (-.68,-.55)--(-.68,1.46) node[above,font=\scriptsize] {$e_2$};
\node[font=\small\bfseries] at (.45,1.58) {20 步后 $E_{20}$};
\fill[blue] (1.2145,.1364) circle (2.5pt); \node[font=\scriptsize,above] at (1.2,.15) {春天};
\fill[teal] (-.4298,.881) circle (2.5pt); \node[font=\scriptsize,right] at (-.43,.88) {花};
\fill[orange] (.1734,-.2188) circle (2.5pt); \node[font=\scriptsize,right] at (.17,-.22) {开了};
\end{scope}
\end{tikzpicture}
\end{center}
\begin{center}\captionof{figure}{训练前后 $E$ 的前两坐标投影：同尺度、非 PCA；不展示第三坐标，不代表真实词义}\label{fig:embedding-before-after}\end{center}
\tagstructend


图中仅取三维向量的前两个坐标画散点（不是 PCA）；
每个点第三坐标仍存在。比较同一 Token 在训练前后的移动，
是为了看到参数会学，并不是用点间距离直接判定词义。

\subsection{Word2Vec 与上下文化表示有什么关系？}

经典 Word2Vec 用词与上下文之间的预测关系训练\textbf{静态词向量}，
体现“上下文相关的词可能学到某些类似特征”的直觉。
现代 Transformer 的 Embedding 通常与后面网络一起训练，
同一个 Token 在不同上下文中的后续隐藏表示还会不同，
所以 Embedding 表里的一行\textbf{不等于最终上下文化含义}。

\begin{warningbox}
本节 $5\times3$ 的 $E$ 是新版统一贯通实验的数据来源；
后续若仍出现 $d=2$ 的 Attention 或 Block 示例，
它们是保留的\textbf{独立局部演示}，不能从这里直接沿用其权重。
其余二维算例将在 M3 与完整 Transformer 的数值链一并迁移。
\end{warningbox}

\par\medskip\noindent{\sffamily\footnotesize\color{blue}\textbf{读完自测}}\quad
\phantomsection\label{check:20-embeddings}
词表大小为 $5$、向量维度为 $3$，查三个 Token ID 得到多大的矩阵？
\quad{\footnotesize\hyperref[answer:20-embeddings]{核对解答。}}\par\medskip

\begin{conclusion}
ID 查表得到可训练特征；梯度会改变这些数字，供后续语言预测使用。
下一章规定模型要学的任务：给什么输入，输出哪一个 Token 的分数？
\end{conclusion}
